\section{数值模拟与结果}
\label{sec:numerics}

在本节中我们给出格点 QCD 结果。我们先介绍格点设置，随后给出
由两点关联得到的准 TMDWFs 结果。
Wilson 圈的结果在 III.C 小节讨论。III.D 小节研究算符混合效应。关于 CS 核的主要结果在 E 小节给出。最后一个 E 小节包含一些总体性的讨论。

\subsection{格点设置 }

我们的数值模拟采用 $N_f=2+1+1$ 的价 clover 费米子置于高度改进 staggered 夸克（HISQ）海~\cite{Follana:2006rc}之上并配以一圈 Symanzik 改进规范作用量~\cite{Symanzik:1983dc}，规范系综由 MILC 合作组以周期边界条件生成~\cite{MILC:2012znn}。计算中我们采用单一系综：晶格间距 $a\simeq0.12$~fm、体积 $n_s^3\times n_t=48^3\times64$、物理的海夸克质量。为了提高信噪比，我们把轻价夸克质量调到与奇异夸克相同，即 $m_{\pi}^{\mathrm{sea}}=130$~MeV 与 $m_{\pi}^{\mathrm{val}}=670$~MeV，
这可能产生一些非幺正性效应。另一方面，Collins-Soper 核对夸克质量的依赖很弱，我们可以把价夸克视为类奇异的，即所涉及的强子为 K 介子。

为进一步改善统计信号，我们在规范系综中采用超立方（HYP）光滑化的粗链~\cite{Hasenfratz:2001hp}。为达到 CS 核的大动量极限，我们采用三个不同的强子动量 $P^z=2\pi/n_s\times\{8,10,12\}=\{1.72,~2.15,~2.58\}$~GeV，对应增速因子 $\gamma=\{2.57,~3.21,~3.85\}$。

%%%%%%%%%%%%%%%%%%%%%%%%%%

\subsection{由两点关联函数得到的准 TMDWFs}
\label{sec:quasi-WF_from_2pt}
为计算式~(\ref{eq:quasiWFinmomentumspace})定义的准 TMDWFs，我们生成库仑规范的墙源传播子
\begin{align}
S_{w}\left(x, t, t^{\prime} ; \vec{p}\right)=\sum_{\vec{y}} S\left(t, \vec{x} ; t^{\prime}, \vec{y}\right) e^{i \vec{p} \cdot(\vec{y}-\vec{x})}, \label{eq:wallsource}
\end{align}
其中 $(t^{\prime},\vec{y})$ 与 $(t,\vec{x})$ 分别表示源与汇的时空位置。于是可以构造与式~(\ref{eq:quasiWFinmomentumspace})中准 TMDWFs 相关的两点函数（2pt）：
\begin{align}
C_2^{\pm}&	\left(z,b_{\perp},P^z; p^z,L,t \right)=\frac{1}{n_s^3}\sum_{\vec{x}}\mathrm{tr} e^{i\vec{P}\cdot\vec{x}}\left\langle S_w^{\dagger}\left(\vec{x}_1,t,0;-\vec{p}\right)\right. \nonumber\\
&\;\;\;\;\;\;\;\;\;\; \;\;\;\;\;\;\;\;\;\;    \times \left. \Gamma U_{\sqsupset, \pm L}\left(\vec{x}_1,\vec{x}_2\right) S_w\left(\vec{x}_2,t,0;\vec{p} \right)  \right\rangle
\label{eq:nonlocal2pt}
\end{align}
其中 $\vec{x}_1=\vec{x}+z\hat{n}_z/2+b_{\perp}\hat{n}_{\perp}$，$\vec{x}_2=\vec{x}-z \hat{n}_{z}/2$。夸克动量 $\vec{p}=(\vec{0}_{\perp},p^z)$ 沿 $z$ 方向，两个夸克各携带强子动量的一半，因此强子动量满足 $\vec{P}=2\vec{p}$。反夸克传播子可以由式~(\ref{eq:1loophardkernel})出发、利用 $\gamma_5$ 厄米性 $S_w(x,y)=\gamma_5S_w(y,x)^{\dagger}\gamma_5$ 得到。如上所述，狄拉克结构取为 $\Gamma=\gamma^z\gamma_5$ 与 $\gamma^t\gamma_5$，它们在大 $P^z$ 极限下可投影到领头扭度的光锥贡献。

通过生成具有夸克动量 $p^z=\pm\{4,5,6\}\times2\pi/(n_sa)$ 的墙源传播子，以及按照式~(\ref{eq:stapleshapedUlink})的三段规范链，就可以在格点上构造两点关联函数。借助约化公式，$C_2^{\pm}$ 可以参数化为
 \begin{align}
C_2^{\pm}&\left(z,b_{\perp},P^z; p^z,L,t \right)=\frac{A_w(p^z)A_p}{2E}\tilde{\Phi}^{\pm0}\left(z,b_{\perp},P^z,L\right) \nonumber\\
& \times e^{-Et}\left[1+c_0\left(z,b_{\perp},P^z,L\right)e^{-\Delta Et} \right], \label{eq:C2parametrization}
 \end{align}
其中 $A_w(p^z)$ 是采用库仑规范固定墙源的赝标介子内插场的矩阵元，而 $A_p$ 是点源（汇）对应的矩阵元。这些项以及因子 $E^{-1}=1/\sqrt{m_{\pi}^2+\left(P^z\right)^2}$ 都被同一时间片上的局域两点函数 $C_2^{\pm}\left(0,0,P^z; p^z,0,t \right)$ 相消，于是剩下的基态矩阵元 $\tilde{\Phi}^{\pm0}$ 被归一化。非定域与局域两点函数之比可以参数化为
\begin{align}
&R^{\pm}\left(z,b_{\perp},P^z,L,t \right)=\frac{C_2^{\pm}\left(z,b_{\perp},P^z,L,t \right)}{C_2\left(0,0,P^z,0,t \right)} \nonumber\\
=&\tilde{\Phi}^{\pm0}\left(z,b_{\perp},P^z,L\right)\left[1+c_0\left(z,b_{\perp},P^z,L\right)e^{-\Delta Et} \right].
\label{eq:two-state-fit}
\end{align}
 %}
其中分母中的 $C_2$ 是局域
关联函数。

 %%%%%%%%%%%%%%%%%%%%%%%%%%%%%%%%%%
\begin{figure}
\centering
\includegraphics[width=0.62\textwidth]{t_dep_b2z0_r}
\caption{提取 $\tilde{\Phi}^{\pm0}\left(z,b_{\perp},P^z,L\right)$ 的两态拟合与单态拟合的比较。以 $\{z,b_{\perp},P^z,L\}=\{0a, 2a, 24\pi/n_s, 6a\}$ 为例可以看到：两态拟合适用于 $t\in[2a, 8a]$，而单态拟合适用于 $t\in[5a, 8a]$。两者的拟合结果相互一致，但单态拟合更为保守。}
\label{fig:tdependenceofratio}
\end{figure}
%%%%%%%%%%%%%%%%%%%%%%%%%%%%%%%%%%


在上述参数化中，激发态的贡献被归入 $c_0$ 项，$\Delta E$ 表示基态与第一激发态之间的能隙。随着欧几里得时间的增大，激发态贡献衰减，在大时间处得到的 $R^{\pm}\left(z,b_{\perp},P^z,L,t \right)$ 平台反映基态贡献 $\tilde{\Phi}^{\pm0}$。我们采用两种方法提取 $\tilde{\Phi}^{\pm0}$：即直接使用式~\eqref{eq:two-state-fit}的两态拟合，以及令 $c_0=0$ 的单态拟合。在足够大的欧几里得时间下，图~\ref{fig:tdependenceofratio}针对小 $\{z,b_{\perp}\}$ 情形给出了两种方法的比较。从该图可以看出，单态拟合结果与两态拟合一致，但给出更保守的误差估计。两态拟合适用于 $t\in[2a, 8a]$，单态拟合适用于平台区 $t\in[5a, 8a]$。而在极高精度下，两态拟合会被激发态污染主导，尤其对于大 $\{z, b_\perp\}$ 的情形。因此按当前精度，我们在后续分析中采用来自单态拟合的更保守结果。更多细节见附录~\ref{appsec:C2}。



%%%%%%%%%%%%%%%%%%%%%%%%%%%%%%%%%%
\subsection{Wilson 圈重正化}
\label{subsec:numerical_wl}
%%%%%%%%%%%%%%%%%%%%%%%%%%%%%%%%%%

由两点函数联合拟合提取的未减除准 TMDWF 矩阵元 $\tilde{\Phi}^{\pm0}\left(z,b_{\perp},P^z,a,L\right)$ 含有来自线性发散的因子 $e^{-\delta \bar{m}(2L+b_{\perp})}$、重夸克有效势因子 $e^{-V(b_{\perp})L}$ 以及对数发散 $Z_{O}$：
\begin{align}
\tilde{\Phi}^{\pm0}\left(z,b_{\perp},P^z,a,L\right) \propto e^{-\delta \bar{m}(2L+b_{\perp})}e^{-V(b_{\perp})L}Z_{O}.
\end{align}
其中 $Z_O$ 对晶格间距 $a$ 具有对数依赖。

$e^{-\delta \bar{m}(2L+b_{\perp})}$ 中的线性发散来自 Wilson 线的自能~\cite{Ji:2017oey,Ishikawa:2017faj,Green:2017xeu, Ji:2020brr}，其中 $\delta\bar m$ 包含一个正比于 $1/a$ 的项和一个非微扰的 renormalon 贡献 $m_{0}$：
\begin{align}
\delta\bar m = \frac{m_{-1}(a)}{a} - m_0  \ .
\end{align}
注意线性发散项的指数正比于 Wilson 链的总长度，例如钉扎链的 $2L+b_{\perp}$。由于这一因子，Wilson 圈的数值在小 $a$ 与大 $L$ 时会急剧下降。

重夸克有效势项 $e^{-V(b_{\perp})L}$ 来自钉扎链中沿 $z$ 方向的两条 Wilson 线之间的相互作用。重夸克有效势 $V(b_{\perp})$ 常被用来确定系综的晶格间距。

对数发散 $Z_{O}$ 来自涉及 Wilson 线与轻夸克的顶点。经重正化群方程重求和并匹配到 $\overline{\rm MS}$ 方案的领头阶对数发散为~\cite{Ji:1991pr, LatticePartonCollaborationLPC:2021xdx}
\begin{align}\label{eq:ZO_RS_higher}
Z_{O}(1/a, \mu)=\left(\frac{\ln [1 /(a \Lambda_{\rm QCD}^{\rm latt})]}{\ln [\mu / \Lambda_{\rm QCD}^{\rm MS}]}\right)^{\frac{3 C_{F}}{b_0}},\nonumber\\
\end{align}
其中 $\Lambda_{\rm QCD}^{\rm latt}$ 不同于 $\overline{\rm MS}$ 方案中的值。可以用这两者有效地吸收高阶贡献~\cite{Lepage:1992xa}。

本文采用 Wilson 圈重正化方法~\cite{Chen:2016fxx,Zhang:2017bzy,Musch:2010ka,Green:2017xeu,Zhang:2017zfe}，其中式~(\ref{eq:WilsonLoop})定义的 Wilson 圈 $Z_E$ 包含线性发散与重夸克有效势：
\begin{align}
Z_E\left(2L, b_{\perp}, a\right) \propto e^{-\delta \bar{m}(4L+2b_{\perp})}e^{-V(b_{\perp})2L}.
\end{align}
根据文献~\cite{LatticePartonCollaborationLPC:2021xdx}，Wilson 圈中的 $\delta \bar{m}$ 与强子矩阵元中的相同，因此预期除以 $\sqrt{Z_E\left(2L, b_{\perp}, a\right)}$ 后线性发散即被去除：
\begin{align}
\tilde{\Phi}^{\pm}(z,b_{\perp},P^z,a,L)=\frac{\tilde{\Phi}^{\pm0}\left(z,b_{\perp},P^z,a,L\right)}{\sqrt{Z_E\left(2L, b_{\perp}, a\right)}}.
\end{align}
如图~\ref{fig:WilsonLooprenormalization}所示，减除后的准 TMDWFs 在 $L\ge 0.4$ fm 时趋于常数。于是我们采用如下定义的减除准 TMDWFs
\begin{align}\label{eq:WLRQuasi}
\tilde{\Phi}^{\pm}(z,b_{\perp},P^z, a)=\lim_{L\to\infty}\frac{\tilde{\Phi}^{\pm0}\left(z,b_{\perp},P^z,a,L\right)}{\sqrt{Z_E\left(2L, b_{\perp}, a\right)}}.
\end{align}
然而，预期仍存在残余的对数发散 $Z_{O}$：
\begin{align}
\frac{\tilde{\Phi}^{\pm0}\left(z,b_{\perp},P^z,a,L\right)}{\sqrt{Z_E\left(2L, b_{\perp}, a\right)}}
\propto \frac{e^{-\delta \bar{m}(2L+b_{\perp})}e^{-V(b_{\perp})L}Z_{O}}{e^{-\delta \bar{m}(2L+b_{\perp})}e^{-V(b_{\perp})L}}
= Z_{O}.
\end{align}
在提取 CS 核时采用的是准 TMDWFs 之比，因而残余对数发散 $Z_{O}$ 被相消。



 %%%%%%%%%%%%%%%%%%%%%%%%%%%%%%%%%%
\begin{figure}
\centering
\includegraphics[width=0.62\textwidth]{l_dep_8_21_r}
\includegraphics[width=0.62\textwidth]{l_dep_8_21_i}
\caption{未减除与减除准 TMDWFs 的 $L$ 依赖结果：实部（上图）与虚部（下图），以及 Wilson 圈的平方根。以 $\Gamma=\gamma^z\gamma_5$ 且 $\{P^z,b_{\perp}, z\}=\{16\pi/n_s,2a,2a\}$ 的情形为例展示。}
\label{fig:WilsonLooprenormalization}
\end{figure}
%%%%%%%%%%%%%%%%%%%%%%%%%%%%%%%%%%

\subsection{算符中的高扭度效应}
\label{sec:operator}


%%%%%%%%%%%%%%%%%%%%%%%%%%%%%%%%%%%%%%
\begin{figure}
\centering
\includegraphics[width=0.62\textwidth]{qwf_co_b3_r_p12}
\includegraphics[width=0.62\textwidth]{qwf_co_b3_i_p12}
\caption{不同狄拉克结构下准 WF 矩阵元的 $\lambda$ 依赖。此处以 $\{P^z,b_{\perp}\}=\{24\pi/n_s,3a\}$ 为例，两情形间的偏差主要来自幂次修正。更多 $\{P^z,b_{\perp}\}$ 组合的结果见附录~\ref{sec:app_operator_mixing}。}
    \label{fig:qwf_coordinate}
\end{figure}
%%%%%%%%%%%%%%%%%%%%%%%%%%%%%%%%%%%%%%


对于欧几里得格点上的赝标介子，$\Gamma=\gamma^t\gamma_5$ 与 $\gamma^z\gamma_5$ 都能投影到领头扭度的光锥分布振幅，即大 $P^z$ 极限下的 $\gamma^+\gamma_5$。两者之间的差异源于按 $M^2/\left(P^z\right)^2$ 表示的幂次修正。

图~\ref{fig:qwf_coordinate}展示了在 $P^z=24\pi/n_s\simeq2.58$~GeV 处，$\Gamma=\gamma^t\gamma_5$ 与 $\gamma^z\gamma_5$ 两种情形下准 TMDWFs 随 $\lambda= z   P^z$ 依赖的比较。由图可见，两组结果的实部在小 $\lambda$ 区域存在一些差异。预期这些差异随 $P^z$ 增大而减小，且含 $\gamma^t\gamma_5$ 与 $\gamma^z\gamma_5$ 的关联函数将从相反方向逐渐收敛到光锥结果。此外，在光锥坐标中 $\gamma^t$ 与 $\gamma^z$ 可用 $\gamma^+$ 与 $\gamma^-$ 表示为
\begin{align}
\gamma^t\gamma_5=\frac{1}{\sqrt{2}}(\gamma^++\gamma^-)\gamma_5,\nonumber\\
\gamma^z\gamma_5=\frac{1}{\sqrt{2}}(\gamma^+-\gamma^-)\gamma_5,
\end{align}
由于动量沿光锥方向，含 $\gamma^-\gamma_5$ 的算符对应 TMDWFs 的高阶项。因此，有限 $P^z$ 引起的幂次修正有望在这两项的平均中被消除：
\begin{align}
\tilde{\Phi}^{\pm}=\frac{1}{2}\left[\tilde{\Phi}^{\pm}\left(\Gamma=\gamma^t\gamma_5\right) + \tilde{\Phi}^{\pm}\left(\Gamma=\gamma^z\gamma_5\right)\right].
\end{align}
定量分析参见文献~\cite{LatticeParton:2020uhz}的附录。算符混合效应可达 5\% 的量级~\cite{LatticeParton:2020uhz}，远小于下一小节讨论的系统不确定度。

根据我们的数值模拟，坐标空间中的减除准 TMDWFs $\tilde{\Phi}^{\pm}(z,b_{\perp},P^z)$ 作为 $\lambda=zP^z$ 的函数是复值的，如图~\ref{fig:qwf_co}所示。示例是 $P^z=24\pi/n_s$、$b_{\perp}=2a$ 与 $4a$ 时 $\tilde{\Psi}^{\pm}(x,b_{\perp},P^z)$ 的实部与虚部。
为确定动量空间中的准 TMDWFs $\tilde{\Psi}^{\pm}(x,b_{\perp},P^z)$，我们采用傅里叶变换（FT）
\begin{align}
\tilde{\Psi}^{\pm}(x,b_{\perp},P^z)=\frac{1}{2\pi}\sum_{z_{\rm min}}^{z_{\rm max}}e^{ixzP^z}\tilde{\Phi}^{\pm}(z,b_{\perp},P^z).
\label{eq:FT}
\end{align}
由于 $\tilde{\Phi}^{\pm}(z,b_{\perp},P^z)$ 存在虚部，$\tilde{\Psi}^{\pm}(x,b_{\perp},P^z)$ 也带有虚部。以 $P^z=24\pi/n_s$ 以及 $b_{\perp}=2a$ 与 $4a$ 为例，我们得到了 $\tilde{\Psi}^{\pm}(x,b_{\perp},P^z)$ 实部与虚部的动量空间准 TMDWFs，示于图~\ref{fig:qwf_momenta}。我们在 $z_{\rm min}$ 与 $z_{\rm max}$ 处截断 FT。$\tilde{\Phi}^{\pm}(z_{\rm min},b_{\perp},P^z)$ 与 $\tilde{\Phi}^{\pm}(z_{\rm max},b_{\perp},P^z)$ 偏离零的程度度量了由此产生的截断误差。对我们数值上能够实现的最大 $z$ 范围——$z_{\rm min}=-1.44$fm、$z_{\rm max}=1.44$fm——该误差仍然可观。这种强行截断的 FT 导致 TMDWFs 出现振荡行为。若在傅里叶变换之前，先对作为 $zP^z$ 函数的 $\tilde{\Phi}(z,b_{\perp},P^z)$ 作适当外推，则可以消除 $\tilde{\Phi}(x,b_{\perp},P^z)$ 中的振荡。但由于我们数据的信噪比不够平滑，最终采用了强行傅里叶变换。

%%%%%%%%%%%%%%%%%%%%%%%%%%%%%%%%%%%%%%
\begin{figure}
\centering
\includegraphics[width=0.62\textwidth]{qwf_co_b2}
\includegraphics[width=0.62\textwidth]{qwf_co_b4}
\includegraphics[width=0.62\textwidth]{qwf_co_i_b2}
\includegraphics[width=0.62\textwidth]{qwf_co_i_b4}
\caption{坐标空间中减除准 TMDWFs 的示例。此处取 $P^z=24\pi/n_s$ 的情形。根据我们的结果，减除准 TMDWFs $\tilde{\Psi}^+(z,b_{\perp},P^z)$ 是复值的，因此其实部与虚部都需要考察。上方四幅图是对两种狄拉克结构（$\gamma^t\gamma_5$ 与 $\gamma^z\gamma_5$）取平均后、作为 $\lambda=zP^z$ 函数的减除准 TMDWFs $\tilde{\Psi}^+(z,b_{\perp},P^z)$。左上图给出 $b_{\perp}=2a$ 时 $\tilde{\Psi}^+(z,b_{\perp},P^z)$ 的实部，右上图则是 $b_{\perp}=4a$ 时的 $\tilde{\Psi}^+(z,b_{\perp},P^z)$。下方两幅图为对应的 $b_{\perp}=2a$ 与 $4a$ 的虚部。}
    \label{fig:qwf_co}
\end{figure}
%%%%%%%%%%%%%%%%%%%%%%%%%%%%%%%%%%%%%%


%%%%%%%%%%%%%%%%%%%%%%%%%%%%%%%%%%%%%%
\begin{figure}
\centering
\includegraphics[width=0.62\textwidth]{qwf_mom_b2}
\includegraphics[width=0.62\textwidth]{qwf_mom_b4}
\includegraphics[width=0.62\textwidth]{qwf_mom_i_b2}
\includegraphics[width=0.62\textwidth]{qwf_mom_i_b4}
\caption{强子动量 $P^z=24\pi/n_s$ 下动量空间减除准 TMDWFs 的示例。动量空间的减除准 TMDWFs $\tilde{\Psi}^+(x,b_{\perp},P^z)$ 由 $\tilde{\Phi}^+(z,b_{\perp},P^z)$ 经傅里叶变换得到，后者具有实部与虚部，两者都必须考察。如式~(\ref{eq:FT})所示，采用强行 FT 确定 $\tilde{\Psi}^+(x,b_{\perp},P^z)$，其中 $z_{\rm min}=-1.44$~fm、$z_{\rm max}=1.44$~fm。左上图给出 $b_{\perp}=2a$ 情形 $\tilde{\Psi}^+(x,b_{\perp},P^z)$ 的实部，右上图则是 $b_{\perp}=4a$ 情形。相应地，下方两幅图为 $b_{\perp}=2a$ 与 $4a$ 时 $\tilde{\Psi}^+(x,b_{\perp},P^z)$ 的虚部。}
    \label{fig:qwf_momenta}
\end{figure}
%%%%%%%%%%%%%%%%%%%%%%%%%%%%%%%%%%%%%%




\subsection{由准 TMDWFs 确定 Collins-Soper 核}
\label{subsec:csresults}

CS 核支配快度演化，因而不依赖于所涉部分子的动量分数。但如式~(\ref{eq:matching})所指出的，该因子化公式仅在 $xP^z\gg \Lambda_{\rm QCD}$ 时成立，在端点区域 $x\to0,1$ 可能失效。幂次修正预计具有 $1/\left(xP^z\right)^2$ 或 $1/\left(\bar{x}P^z\right)^2$ 的形式。因此，数值 CS 核被拟合为 $x$、$P_1^z$ 与 $P_2^z$ 的函数，记作 $K(b_{\perp},\mu,x,P^z_1,P^z_2)$，
\begin{align}
&K(b_{\perp},\mu,x,P^z_1,P^z_2) =\nonumber\\
&\frac{1}{2\ln(P_1^z/P_2^z)} \left[\ln\frac{H^{+}(xP_2^z,\mu)\tilde{\Psi}^{+}(x,b_{\perp},\mu,P_1^z)}{H^{+}(xP_1^z,\mu)\tilde{\Psi}^{+}(x,b_{\perp},\mu,P_2^z)}\right.  \nonumber\\
	&+\left.\ln\frac{H^{-}(xP_2^z,\mu)\tilde{\Psi}^{-}(x,b_{\perp},\mu,P_1^z)}{H^{-}(xP_1^z,\mu)\tilde{\Psi}^{-}(x,b_{\perp},\mu,P_2^z)}\right].
\end{align}
这里 $K(b_{\perp},\mu,x,P^z_1,P^z_2)$ 由一圈匹配的微扰匹配 kernel 与准 TMDWFs 提取，将带有形如 $\mathcal{O}\left(1/\left(xP^z\right)^2\right)$ 与 $\mathcal{O}\left(1/\left(\bar{x}P^z\right)^2\right)$ 的幂次修正。为提取领头幂次贡献，我们采用如下参数化
\begin{align}
&K(b_{\perp},\mu,x,P^z_1,P^z_2) = K(b_{\perp},\mu) \nonumber\\
&\quad+A\left[\frac{1}{x^2(1-x)^2(P^z_1)^2}-\frac{1}{x^2(1-x)^2(P^z_2)^2}\right], \label{eq:parametrizationofCSkernel}
\end{align}
其中 $A$ 是计及领头高幂次贡献的系数，可以通过在不太靠近 $x=0, 1$ 的区域对不同格点数据做联合拟合来确定。

图~\ref{fig:K_b}给出了 $b_{\perp}=\{0.12,0.24,0.36,0.48,0.60\}$fm 的物理 CS 核 $K(b_{\perp},\mu)$。利用 $P^z=\{8,10,12\}\times2\pi/n_s$ 三种准 TMDWFs，可以提取 $P_1^z/P_2^z=10/8$ 与 $12/8$ 的 $K(b_{\perp},\mu,x,P^z_1,P^z_2)$，以不同色带表示。除端点区域（$x<0.2$ 或 $x>0.8$）之外，格点数据是平坦的，反映了领头幂次贡献，符合预期。利用参数化公式式~(\ref{eq:parametrizationofCSkernel})，通过对数据拟合可以确定物理 CS 核 $K(b_{\perp},\mu)$，以绿色带表示。如前所述，在大 $b_{\perp}$ 处，准 TMDWFs 因傅里叶变换的截断而出现振荡，这也会影响所提取的 $K(b_{\perp},\mu,x,P^z_1,P^z_2)$，见图~\ref{fig:K_b}下图所示。原则上，一旦能得到更大 $z$ 的数据，或者知道如何把当前数据外推到更大的 $z$ 或 $\lambda$ 区域，这种振荡效应即可消除。


%%%%%%%%%%%%%%%%%%%%%%%%%%%%%%%%%%%%%%
\begin{figure}
\centering
    \includegraphics[width=0.62\textwidth]{K_b1}
    \includegraphics[width=0.62\textwidth]{K_b2}
    \includegraphics[width=0.62\textwidth]{K_b3}
    \includegraphics[width=0.62\textwidth]{K_b4}
    \includegraphics[width=0.62\textwidth]{K_b5}
\caption{由准 WFs $\tilde{\Psi}^{\pm}$ 提取的 $K(b_{\perp},\mu,x,P^z_1,P^z_2)$ 拟合结果。所选动量对 $\{P_1^z,P^z_2\}$ 在图例中以 $P^z_1/P_2^z$ 标注。各图对应 $b_{\perp}=\{0.12,0.24,0.36,0.48,0.60\}$fm 的情形。} 水平阴影带给出 $K(b_{\perp},\mu)$ 的中心值与不确定度以及文中所述的 $x$ 拟合范围。在大 $b_{\perp}$ 区域，两端阴影区（LaMET 方法失效之处）存在强振荡，它由大动量展开的失效造成。
    \label{fig:K_b}
\end{figure}
%%%%%%%%%%%%%%%%%%%%%%%%%%%%%%%%%%%%%%


理论上物理 CS 核纯为实数，但在一圈匹配下仍存在残余虚部。如上所述，这个虚部项来自微扰匹配 kernel。容易证明 $H^{\pm}(zP_2^z,\mu)/H^{\pm}(zP_1^z,\mu)$ 含有虚部，而格点结果 $\tilde{\Psi}^{\pm}(x,b_{\perp},\mu,P_1^z)/\tilde{\Psi}^{\pm}(x,b_{\perp},\mu,P_2^z)$ 近乎为实数，见 \ref{fig:ratio_phi}。因此，我们把该虚部视为来自因子化定理的系统不确定度。此不确定度可表示为

\begin{align}
&\sigma_{\mathrm{sys}}=\sqrt{K(b_{\perp},\mu)+\text{Im}^2\left[K^+(b_{\perp},\mu)\right]}-K(b_{\perp},\mu)
\label{eq:systematical_error}
\end{align}
其中 $\text{Im}\left[K^+(b_{\perp},\mu)\right]$ 表示仅用 $\tilde{\Psi}^+$ 提取 $K(b_{\perp},\mu)$ 所得数值结果的虚部：
\begin{align}
K^+(b_{\perp},\mu)=&\frac{1}{\ln (P^z_1/P^z_2)}\ln\frac{H^+(xP^z_2,\mu)\tilde \Psi^+(x, b_\perp, \mu, P^z_1)}{H^+(xP^z_1,\mu)\tilde \Psi^+(x, b_\perp, \mu, P^z_2)}.
\label{eq:K+}
\end{align}

应当注意，微扰匹配 kernel $H^+(zP^z,\mu)$ 是 $H^-(zP^z,\mu)$ 的复共轭，也就是说当我们取 $H^+(zP_2^z,\mu)/H^+(zP_1^z,\mu)$ 与 $H^-(zP_2^z,\mu)/H^-(zP_1^z,\mu)$ 的平均时，其中的虚部可以相互抵消。因此，作为最终结果我们取 $K(b_{\perp},\mu)=\left[K^+(b_{\perp},\mu)+K^-(b_{\perp},\mu)\right]/2$ 以保留实部，并把虚部贡献视为我们的系统不确定度。

%%%%%%%%%%%%%%%%%%%%%%%%%%%%%%%%%%%%%%
\begin{figure}
\centering
\includegraphics[width=0.62\textwidth]{imag_ratio_b3}
\caption{$b_{\perp}=3$a 时不同 $P^z$ 下准 TMDWFs 之比 $\frac{\tilde{\Psi}^+(x,b_{\perp},\mu,P_1^z)}{\tilde{\Psi}^+(x,b_{\perp},\mu,P_2^z)}$ 的虚部作为动量分数 $x$ 函数的数值结果示例。两种情形的虚部都接近零。}
    \label{fig:ratio_phi}
\end{figure}
%%%%%%%%%%%%%%%%%%%%%%%%%%%%%%%%%%%%%%

%%%%%%%%%%%%%%%%%%%%%%%%%%%%%%%%%%%%%%

\subsection{结果与讨论}

应当注意，格点上经 Wilson 圈重正化的准 TMDWF（式~(\ref{eq:WLRQuasi})）带有对标度 $a$ 的依赖。如果通过除以 $Z_{O}$（式~(\ref{eq:ZO_RS_higher})）把它转换到 $\overline{\rm MS}$ 方案，就会引入标度 $\mu$。原则上应当把 Wilson 圈重正化的准 TMDWF 转换到 $\overline{\rm MS}$ 方案，因为我们的因子化公式在该方案下成立。然而，由于 $Z_{O}$ 不依赖动量 $P_{z}$，它在准 TMDWFs 之比中相消，标度依赖也随之消去。因此，在提取 CS 核的过程中无需对准 TMDWF 做方案与标度转换。

由不同动量准 TMDWFs 之比联合拟合提取的 CS 核以红色数据点示于图~\ref{fig:K_com}。{图中我们对 $K(b_{\perp},\mu)$ 给出两类误差：较小的是统计不确定度，较大者同时包含统计与系统不确定度}。在小 $b_{\perp}$ 区域，由于大的幂次修正与非零虚部，系统不确定度占主导。

作为比较，我们还给出 CS 核的树图级匹配结果。取领头阶匹配 kernel $H(xP^z,\mu)=1+\mathcal{O}(\alpha_s)$，式~(\ref{eq:extractingCSkernel})简化为 $z=0$ 处、动量比 $P_1^z/P^z_2$ 下准 TMDWFs $\tilde{\Phi}$ 之比。图~\ref{fig:K_com}中的蓝点表示树图级匹配所得结果，仅显示统计不确定度。

%%%%%%%%%%%%%%%%%%%%%%%%%%%%%%%%%%%%%%
\begin{figure}
\centering
\includegraphics[width=0.62\textwidth]{K_com}
\includegraphics[width=0.62\textwidth]{K_pheno}
\caption{{上图给出我们的结果 $K(b_{\perp},\mu)$ 与 $K_0(b_{\perp},\mu)$ 同 SWZ~\cite{Shanahan:2021tst}、LPC~\cite{LatticeParton:2020uhz}、ETMC/PKU~\cite{Li:2021wvl} 与 SVZES~\cite{Schlemmer:2021aij} 各格点结果以及最高三圈微扰计算的比较。$K(b_{\perp},\mu)$ 表示经一圈匹配提取的 CS 核，其不确定度包括统计误差以及来自非零虚部的系统误差。$K_0(b_{\perp},\mu)$ 表示我们的树图级结果，仅带统计不确定度。下图给出我们的结果与唯象抽取的比较：SV19~\cite{Scimemi:2019cmh}、Pavia19~\cite{Bacchetta:2019sam} 与 SIYY15~\cite{Sun:2014dqm} 给出拟合数据得到的 CS 核唯象参数化，所用数据来自} Drell-Yan 之类的高能碰撞过程。}
\label{fig:K_com}
\end{figure}
%%%%%%%%%%%%%%%%%%%%%%%%%%%%%%%%%%%%%%

我们将自己的结果与微扰计算、唯象抽取以及其他合作组确定的格点结果进行比较。

图~\ref{fig:K_com}上图中的黑色实线与虚线表示最高三圈的微扰结果，跑动耦合常数取 $\alpha_s(\mu=1/b_{\perp})$。微扰计算在小 $b_{\perp}$ 区域（$b_{\perp}\ll1/\Lambda_{\mathrm{QCD}}$）工作良好，但随着 $b_{\perp}$ 增大会发散。相比之下，格点计算将在非微扰区给出精确预言，而由于幂次修正，
它在小 $b_{\perp}$ 区域可能遭受较大的系统不确定度。

与本工作类似，LPC~\cite{LatticeParton:2020uhz} 与 ETMC/PKU~\cite{Li:2021wvl} 的结果也是经由树图级匹配从准 TMDWFs 中提取的。采用一圈匹配公式并考虑算符混合效应，将有助于减小系统不确定度并获得更精确的结果。此外，考虑规范链的不同方向有助于消除非物理虚部的贡献，从而提高结果的精度。

另一方面，SWZ~\cite{Shanahan:2021tst} 与 SVZES~\cite{Schlemmer:2021aij} 的结果是从准 TMDPDFs 得到的。与复杂的核子关联函数相比，介子的关联函数更容易获得较好的信号。另外，介子波函数在 $x$ 空间近似对称，因而更便于参数化振荡效应并在大 $P^z$ 极限下获得物理结果，如图~(\ref{fig:K_b})那样。再者，轻介子更容易达到较大的增速因子；从小 $b_{\perp}$ 区域可以看到，来自准 TMDWFs 的结果比来自准 TMDPDFs 的结果与微扰计算吻合得更好。

图~(\ref{fig:K_com})的下图给出与唯象结果的比较。SV19~\cite{Scimemi:2019cmh} 与 SIYY15~\cite{Sun:2014dqm} 采用包含微扰与非微扰部分的参数化。而 Pavia19~\cite{Bacchetta:2019sam} 则借助 TMDPDFs 的因子化、由快度导数获得 CS 核；此外他们还从 Drell-Yan 数据拟合参数以得到其唯象 CS 核。不同方法的结果在非微扰区域表现出明显的不一致。我们的结果显示出与 SV19 更好的一致性。
